\scanpage{198}\mspage{147}
l'intégrale étant l'intégrale \uline{lisse}. (Il y a un énoncé dual concernant l'intégrale propre de

\diagram{\begin{tikzcd}[ampersand replacement=\&,row sep=0pt,column sep=1.2em,cramped] \& X_1 \\ X_0\arrow[ur]\arrow[dr] \& \\ \& X_2 \end{tikzcd}\(\,.\))}

\deleted{Pour ce} Pour que $\varphi_Q$ soit \inserted{$W$-}lisse, il faut et s. que $f_1,f_2,\unsure{f_0}=f_1\alpha_1=f_2\alpha_2$ le soient. (\uline{Moyennant $W(7\text{bis})$})

% slanted
\margin{Pour la partie directe ($f_0,f_1,f_2$ lisses $\Longrightarrow\varphi_Q$ lisse), on n'a besoin que de $W(7\text{bis})$ forme faible. Ce n'est que pour la réciproque (qu'on n'aura jamais à utiliser) qu'on a besoin de $W(7\text{bis})$ (et dans tte sa force).}

% slanted
\margin{On n'a pas $=$ besoin de $W(7\text{bis})$ \inserted{faible}, cf annotation page 156. On utilise une ``\unsure{hypothèse}'' \unsure{de lissité|de rétraction} sur \unsure{$W$|les} \unsure{pour} la \unsure{$\Longrightarrow$} \unsure{réciproque} seulement ($\varphi$ lisse $\Longrightarrow f_0,f_1,f_2$ lisses).}

La fibre de $\varphi_Q$ en un \deleted{\unsure{pt}} élément \unsure{$y$} est isomorphe à

\diagram{\(\text{(**)}\qquad\deleted{\unsure{y}}\qquad \int \begin{tikzcd}[ampersand replacement=\&,row sep=0pt,column sep=1.3em,cramped] \& X_{1y}\\ X_{0y}\arrow[ur]\arrow[dr] \& \\ \& X_{2y}\end{tikzcd}\quad (\simeq\unsure{\mathscr X|\mathcal X}_y).\)}

\uline{Corollaire 1} Supposons $f_1,f_2$ et $f_1\alpha_1=f_2\alpha_2$ $W$-lisses, et les catégories intégrales $(**)$ \deleted{\unsure{de}} \inserted{$W$-}asphériques pour \deleted{les} $y\in\mathrm{Ob}\,Y$. [NB Il s'agit ici de l'intégrale lisse, mais peu importe en fait, car on a vu \deleted{que} (cf \deleted{volet} \deleted{\unsure{gauche}} partie supérieure des diagrammes p. 140, qui impliquent que $\mathscr X$ $W$-asph. \unsure{sss|ssi} $\widetilde X$ l'est) que l'une est $W$-asph. sss l'autre l'est.]

% slanted
\margin{Valable sans hypothèse sur $W$ (autre que Loc(1,2,\unsure{3bis})).}

Alors $\varphi_Q\in W$, i.e. le carré $\unsure{(\alpha)|(*)}$ est $W$-homotopiquement \deleted{trivial}\inserted{cartésien}. Par dualité, $=$ conclusion si on suppose $f_1,f_2,f_1\alpha_1=f_2\alpha_2$ \uline{propres} au lieu de lisses. (car $Q$ $W$-homotopiquement trivial, sss $Q^\circ$ l'est, par prop. 14.)
% archive sheet number: 198
